\section{Introduction}
\label{sec:introduction}

Novel view synthesis (NVS) has witnessed remarkable progress with 3DGS~\cite{kerbl20233d}, which represents scenes as 
collections of anisotropic 3D Gaussians to achieve photorealistic rendering 
at real-time speeds. 
While 3DGS achieves real-time rendering and high visual quality on densely 
captured scenes, its performance degrades dramatically under 
sparse-view settings commonly encountered in practical scenarios.

% 这有点长了
The degradation stems from a fundamental issue in 3DGS optimization: 
\textbf{lack of intrinsic consistency between geometry and appearance}.
In 3DGS, each primitive is parameterized by coupled intrinsic properties, 
including geometric attributes (position $\boldsymbol{\mu}$, covariance 
$\boldsymbol{\Sigma}$, opacity $\alpha$) defining its spatial structure, 
and appearance attributes (view-dependent color $\mathbf{c}(\mathbf{d})$) 
determining its photometric contribution.
For faithful scene reconstruction, these properties must satisfy 
\emph{intrinsic consistency}: geometry should accurately capture the underlying 
3D structure, while appearance should coherently reflect surface photometry 
across viewpoints.
However, sparse-view 3DGS violates this consistency. 
The standard 3DGS optimization relies on per-view photometric supervision that 
independently minimizes rendering loss for each training view.
With limited observations, this independent supervision allows appearance to 
overfit individual views by compensating for geometric errors, while 3D geometry 
remains severely underconstrained due to lack of explicit multi-view regularization.
This leads to internally inconsistent 
Gaussians that produce plausible renderings on training views but severe 
artifacts (floaters or blurriness) on novel views.
These challenges are particularly pronounced in weakly-textured regions, 
where the absence of distinctive appearance cues further exacerbates geometric 
ambiguity.


Addressing this requires jointly solving two coupled challenges. 
\textbf{First}, how to effectively constrain geometry under sparse observations? 
Recent studies have employed pretrained depth estimation models~\cite{zhu2024fsgs,li2024dngaussian,xiong2023sparsegs, xu2025depthsplat} 
to regularize 3DGS geometry, but such approaches suffer from scale ambiguity and 
noise introduced by imperfect pretrained models. 
Other methods~\cite{zhu2024fsgs,zheng2025nexusgs} rely on dense initialization, 
yet the initial geometric cues tend to be gradually forgotten during subsequent optimization.
\textbf{Second}, how to couple accurate geometry with reliable appearance optimization to prevent overfitting? 
BinocularGS~\cite{han2024binocular} builds virtual binocular pairs from rendered depth 
to enforce disparity consistency, but this approach relies on rendered depth whose 
reliability is not guaranteed, potentially propagating depth errors into appearance optimization 
and limiting rendering quality.


To this end, we propose \textbf{ICO-GS} (\textbf{I}ntrinsic Geometry-Appearance 
\textbf{C}onsistency \textbf{O}ptimization for Sparse-view \textbf{G}aussian 
\textbf{S}platting), a principled framework that restores intrinsic consistency 
in sparse-view 3DGS through synergistic geometry-appearance optimization.
Our key insight is that faithful geometry and appearance emerge from their 
\emph{mutual reinforcement}: well-constrained geometry guides appearance to 
learn view-consistent photometry, while reliable appearance supervision in 
turn refines geometry.
We realize this synergy through two coupled components.


\textbf{Geometric regularization via multi-view photometric consistency.}
Under ideal conditions, a 3D point observed from multiple viewpoints should 
exhibit photometric consistency. However, illumination variations, occlusions, 
and monocularly visible regions violate this assumption. 
We therefore adopt feature-based multi-view consistency to regularize geometry, 
mitigating the impact of lighting and imaging variations. 
To handle occlusions prevalent in sparse-view settings, we employ pixel-wise 
top-$k$ selection: for each pixel, we compute photometric errors across all 
source views and retain only the $k$ most consistent ones, robustly filtering 
occluded or unreliable observations. 
For monocularly visible regions, we further incorporate an edge-aware depth 
smoothness term that enforces local coherence while preserving sharp geometric 
boundaries. 
These complementary constraints yield geometry that is both multi-view consistent 
and structurally sound, proving particularly effective in weakly-textured 
regions where photometric cues are scarce.

\textbf{Geometry-guided appearance optimization via virtual view consistency.}
To couple geometry with appearance, we leverage the regularized depth to 
synthesize virtual views, propagating geometric constraints to appearance.
However, naively using all depth estimates introduces noise from uncertain regions.
We address this via cycle-consistency depth filtering: we project each pixel's 
depth to source views and back-project to the original view, retaining only 
pixels with consistent spatial locations.
The filtered depth then guides virtual view synthesis through depth-conditioned 
warping, providing geometric supervision that encourages appearance to capture 
view-consistent photometry rather than overfit individual observations.


Our contributions are summarized as follows:
\begin{itemize}[leftmargin=*,nosep]
      \item We identify \emph{intrinsic consistency}, defined as the coupled correctness 
      of geometry and appearance, as a fundamental principle for sparse-view 3DGS, 
      and reveal how its violation causes severe degradation in novel views.
    \item We address geometry underconstraint through feature-based multi-view 
          photometric regularization with pixel-wise top-$k$ consistency and 
          edge-preserving smoothness.
    \item We prevent appearance overfitting through geometry-guided optimization 
          that synthesizes virtual views from cycle-filtered depth, coupling 
          geometric accuracy with photometric quality.
    \item Extensive experiments validate that ICO-GS achieves state-of-the-art sparse-view novel view synthesis across diverse scenes, particularly on weakly-textured data, with 1.1 dB PSNR gains over prior arts on 3-view DTU scenarios.
\end{itemize}


